\documentclass[conference]{IEEEtran}
\usepackage[utf8]{inputenc}
\usepackage[T1]{fontenc}
\usepackage{cite}
\usepackage{amsmath,amssymb}
\usepackage{graphicx}
\usepackage{booktabs}
\usepackage{url}

\title{Verifier‑Assisted versus LLM‑Only Pre‑deployment Network‑Change Validation: A Preregistered Paired Study}

\author{
\IEEEauthorblockN{Autonomous AI Researcher}
\IEEEauthorblockA{CNSM 2026 Agentic AI Researcher Track}
}

\begin{document}

\maketitle

\begin{abstract}
We measured whether integrating a deterministic pre‑deployment verifier with a bounded automated repair budget increases the fraction of intent‑driven network‑change proposals that pass semantic validation compared with accepting a single LLM candidate. In a preregistered paired design (study id cnsmp2026-llm-verifier-predeployment-001) we executed 300 paired intents (100 per difficulty stratum: low/medium/high) using gpt-5-mini and deterministic\_netops\_validator\_v5. Each paired intent produced one shared\_initial generation that was evaluated under two conditions: baseline (accept single shot) and guarded (deterministic validation and <=1 automated repair). Outcomes: baseline 299/300 unflagged passes (0.9966667), guarded 300/300 (1.0); contingency n11=299, n10=1, n01=0, n00=0; absolute paired difference = 0.0033333; exact McNemar two‑sided p = 1.0; 95\% paired bootstrap percentile CI = [0.00,0.01] (bootstrap\_seed = 1729, B = 10000). Under the supplied synthetic suite and repair budget the single‑shot generator was near ceiling, limiting observable deterministic rescue. We archive preregistration, execution, and analysis artifacts and interpret results with respect to estimand conditioning, validator coverage, repair budget, and operational deployment policies.
\end{abstract}

\section{Introduction}
Generative large language models (LLMs) are being applied to network operations (NetOps) tasks such as intent\ensuremath{\rightarrow}configuration translation, configuration synthesis, and automated repair. For pre‑deployment network changes, semantic correctness properties---sequencing constraints, reachability, and policy preservation---are operationally critical and cannot be assured by superficial syntactic checks alone. A common engineering pattern is to pair generative models with deterministic validators and bounded automated repair (generate\ensuremath{\rightarrow}validate\ensuremath{\rightarrow}repair) to reduce operational risk. However, the measurable benefit of such guarded pipelines depends on the experimental estimand (what is held fixed), the generator single‑shot quality, the validator's expressiveness, and the permitted repair budget.

We preregistered a paired experiment (study id cnsmp2026-llm-verifier-predeployment-001) to quantify the deterministic rescue achievable when the same single LLM candidate is evaluated under two treatments. For each intent the pipeline produced exactly one shared\_initial candidate using the declared generator (gpt-5-mini) and compared: (a) baseline --- accept the single candidate without repair, and (b) guarded --- run a deterministic validator (deterministic\_netops\_validator\_v5) and allow at most one automated repair which itself is re‑validated. The primary research question: does adding a deterministic verifier plus at most one automated repair increase the probability that a single LLM candidate passes semantic pre‑deployment validation compared with accepting that same candidate without repair?

Contributions of this artifact‑grounded study: (1) a preregistered paired evaluation isolating deterministic rescue under a shared\_initial estimand that conditions on a single LLM draw per intent; (2) a complete execution and analysis bundle including validator traces and the single automated repair transcript archived with the run; and (3) an evidence‑grounded interpretation of estimator conditioning, ceiling effects, and operational implications for NetOps pipelines.

\section{Related Work}
Prior work on LLMs for network management has pursued two complementary directions that frame our study. First, multiple recent surveys and systems reports document active efforts to apply LLMs to NetOps tasks including intent translation, template synthesis, protocol test generation, and zero‑touch workflows [https://openalex.org/W4402742293; https://openalex.org/W4411015066; https://openalex.org/W4415275309]. These surveys synthesize a range of prototypes and experimental systems showing that LLMs can produce syntactically plausible device configurations and candidate workflows, but that practical deployment depends critically on downstream validation and environment modeling. For example, PreConfig (a pretrained model for automating configurations) and other synthesis systems demonstrate automation potential while also highlighting dataset and device‑scope differences that limit straightforward generalization [https://openalex.org/W4392875514; https://openalex.org/W4411015066].

Second, a line of systems research specifically proposes and evaluates tool‑augmented LLM architectures that interpose external checkers, planners, or verifiers. Planner‑centric and ReAct‑style frameworks explicitly enable LLMs to call validators or external tools to check semantics and coordinate multi‑step changes; these architectures aim to reduce hallucination and logical errors by shifting some checks outside the model [doi:10.1609/aaai.v40i40.40676; https://openalex.org/W4389518608]. Empirical system reports in the NetOps space have implemented generate\ensuremath{\rightarrow}validate\ensuremath{\rightarrow}repair patterns and reported task‑dependent improvements in correctness when validators capture operational invariants: MeshAgent and related systems show that tool augmentation can improve reliability for certain configuration tasks, while program‑repair‑oriented work demonstrates that combining symbolic fault localization with generative repair has practical value on constrained benchmarks [https://openalex.org/W4416927413; https://openalex.org/W4404240614].

Work on error detection and mitigation complements these architecture proposals. Detection methods grounded in model uncertainty or semantic measures (for example, entropy‑based detectors) and deterministic safety‑net designs have been proposed to catch hallucination or logical inconsistencies before deployment; empirical evaluations indicate these detectors can reduce false assurances but are sensitive to detector design and the class of errors targeted [https://openalex.org/W4399803256; doi:10.2139/ssrn.6538460]. Similarly, domain‑targeted testing work has shown that LLM agents can generate protocol test cases or device‑specific actions, but that correctness requires encoding topological and sequencing invariants into validators or test harnesses to catch semantic failures [https://openalex.org/W4415275309].

Two methodological gaps in the literature motivate our paired shared\_initial design. Many prior evaluations conflate gains from resampling (drawing multiple independent LLM candidates) with gains that arise solely from deterministic validation and repair; permissive generation policies permit the former and thus measure a different operational tradeoff than a single‑call budget. Second, while some works evaluate generate\ensuremath{\rightarrow}validate\ensuremath{\rightarrow}repair pipelines on task‑specific suites, fewer public benchmarks and systematic comparisons isolate semantic correctness (reachability, sequencing, forbidden‑template preservation) as a first‑class evaluation target across difficulty strata [https://openalex.org/W4411015066; https://openalex.org/W4392875514]. Our study explicitly conditions on a single shared\_initial candidate per intent and uses a deterministic validator with an archived, bounded repair operator to provide an operationally meaningful bound on deterministic rescue that complements existing multi‑call or planner‑guided evaluations.

\section{Methodology}
Preregistration and estimand. The study preregistration (cnsmp2026-llm-verifier-predeployment-001; preregistration SHA‑256 = 425cb165\allowbreak{}2354120c\allowbreak{}f48f686b\allowbreak{}964723d6\allowbreak{}4c76d0a8\allowbreak{}abe53cd5\allowbreak{}b5ad0eaa\allowbreak{}1af71fc7) fixed the primary estimand as the absolute paired difference in the unflagged verifier‑pass proportion (guarded minus baseline) computed on 300 paired cases after applying inclusion/exclusion rules. The confirmatory analysis executor paired\_binary\_analysis\_v1 and all inclusion/exclusion rules were frozen prior to the main run. An optional pilot (n = 60) was described in the preregistration for discordance estimation; that pilot was not executed and no post‑preregistration power recalculation was performed (this deviation is recorded in the archived preregistration and Disclosure Statement).

Task generation and synthetic suite. Tasks were produced deterministically by the registered deterministic\_netops\_task\_generator\_v5 into a synthetic\_topology\_suite stratified into low/medium/high difficulty bins by generator metadata (changed\_interface\_count, dependency\_rule\_count, required\_command\_count). Each task manifest entry encodes the full intent, initial and expected final states, required\_sequence constraints, preserve\_unspecified\_state invariants, and protected\_interfaces. Representative difficulty patterns include 'shutdown\_configure\_restore' (medium) and 'paired\_link\_atomic\_mtu\_change' (hard). The synthetic suite enforces validator‑checkable invariants used for scoring.

Shared\_initial protocol and generator. To implement the shared\_initial estimand the pipeline produced exactly one shared\_initial generation per pair and reused that identical candidate for both baseline and guarded conditions. The declared generator for shared\_initial candidate production was gpt-5-mini (execution/model\_configuration.json recorded requested\_model = "gpt-5-mini", declared\_model\_version = "gpt-5-mini"). This isolates deterministic validator/repair rescue from gains achievable by independent resampling, verifier‑guided generation, or iterative generation.

Validator, scoring, and repair operator. deterministic\_netops\_validator\_v5 executed the preregistered full\_intent\_constraint\_validation rule: deterministic checks for sequencing invariants (must\_be\_down\_for\_fields), required\_command\_count, preservation of unspecified state, topology reachability invariants encoded in the synthetic suite, and forbidden‑template protection for protected\_interfaces. Scoring was binary: unflagged == pass. The guarded pipeline permitted at most one automated repair attempt (repair budget <=1). The repair routine was a constrained edit operator limited to localized sequencing or insertion edits; it logged a repair provider trace and was applied only when the validator flagged the shared\_initial candidate. Candidates remaining flagged after the permitted repair counted as failures for the primary estimand.

Contamination controls, negative controls, and stopping rules. The preregistration included one intentionally broken negative control per complexity stratum and deterministic invariant checks (pre/post reachability matrices). Per preregistration, pairs with both calls failing due to provider/infrastructure outages would be excluded (none occurred). The preregistered stopping rule would halt if >20\% dual‑call exclusions occurred; this did not trigger.

Execution instrumentation and audit. The confirmatory run used adapter hosted\_netops\_gvr\_v1 under execution\_mode = scientific\_confirmatory. Execution manifest values (archived): planned\_episode\_count = 600, completed\_episode\_count = 600, complete\_pair\_count = 300, model\_calls\_used = 301, maximum\_model\_calls = 600. Provider call auditing recorded hosted\_model\_call\_event\_count = 301, cache\_miss\_count = 301, stage\_counts: shared\_initial\_generation = 300 and repair = 1. Deterministic reconciliation confirms episode and provider‑call accounting and reports all\_deterministic\_consistency\_checks\_passed = true. These audit values are archived for third‑party verification.

Analysis executor and inferential settings. The preregistered analysis executor paired\_binary\_analysis\_v1 computed paired contingency counts (n11,n10,n01,n00), the absolute paired difference (guarded - baseline), an exact two‑sided McNemar test, and a paired bootstrap percentile 95\% CI. The bootstrap parameters archived with the analysis are seed = 1729 and resamples B = 10000. Sensitivity analyses specified in the preregistration (failure‑as‑zero and stratum‑wise tests with Benjamini--Hochberg FDR control) were run as exploratory artifacts; the confirmatory results reported here follow the primary analysis plan.

\section{Results}
Execution integrity and accounting. The preregistered adapter completed the confirmatory run with completed\_episode\_count = 600, yielding complete\_pair\_count = 300 paired observations under the shared\_initial protocol. Provider/audit records show hosted\_model\_call\_event\_count = 301 (300 shared\_initial generations plus one repair invocation), cache\_miss\_count = 301, and stage\_counts indicating shared\_initial\_generation = 300 and repair = 1. Deterministic reconciliation confirmed consistent episode accounting (all\_deterministic\_consistency\_checks\_passed = true) and the absence of provider/infrastructure failures or excluded pairs.

Primary confirmatory outcomes. The binary scoring produced baseline\_success\_count = 299 and guarded\_success\_count = 300 (analysis/condition\_summary.csv), i.e., baseline = 299/300 (0.9966666667) and guarded = 300/300 (1.0). The paired contingency table (analysis/paired\_contingency\_table.csv) is n11 = 299, n10 = 1, n01 = 0, n00 = 0. The preregistered exact two‑sided McNemar test (operating on the discordant counts n10 and n01) returned p = 1.0. The absolute paired improvement (guarded - baseline) equals 0.0033333333333332993. The paired bootstrap percentile 95\% confidence interval computed with seed = 1729 and B = 10000 is [0.00, 0.01]. Because discordance is minimal (m = n10 + n01 = 1), the exact test has limited sensitivity and the bootstrap CI bounds the plausible absolute improvement to at most one percentage point in this experimental context.

Per‑stratum diagnostics. Per the task manifest stratification, archived per‑stratum tallies were: low difficulty baseline = 100/100 and guarded = 100/100; medium baseline = 100/100 and guarded = 100/100; high baseline = 99/100 and guarded = 100/100. Thus all discordant evidence is concentrated in the high stratum (1 rescues), while low and medium strata were at ceiling under both conditions. These counts are reported in the archived condition\_summary.csv and corroborate that the observed discordance derives from a single hard example under the synthetic generator's difficulty taxonomy.

Repair diagnostic and episode‑level trace. The solitary discordant pair invoked the permitted automated repair. Validator traces (deterministic\_netops\_validator\_v5, validator\_version = 5.0) attribute the initial failure to a sequencing invariant violation (the preregistered must\_be\_down\_for\_fields rule): the shared\_initial candidate omitted an explicit 'admin down' step prior to a VLAN/MTU change required by the intent. The constrained repair operator performed a single localized insertion of the missing sequencing command, after which the deterministic validator reported valid == true for the corrected candidate. The repair therefore produced a deterministic rescue with one additional model/repair stage call. The repair provider trace and the validator's before/after traces for that episode are archived; the manuscript conveys the semantic summary rather than reproducing verbatim provider logs in the body.

Contamination, missingness, and negative controls. The preregistered contamination checks produced an empty contamination\_summary.csv (no contamination flags). The missingness summary reports zero failed or unscorable episodes and zero excluded pairs under preregistered rules (analysis/missingness\_summary.csv). Negative controls included by the preregistration did not cause unexpected pass behavior in the main run; any negative‑control anomalies would have been treated as contamination per the preregistration and are not present in the archived run.

Cost and operational footprint. The guarded pipeline's deterministic rescue required a single additional repair stage call across 300 pairs (model\_calls\_used = 301), indicating that bounded automated repair under the chosen edit operator imposes minimal extra model‑call cost in this run. The empirical repair\_stage\_count = 1 (analysis artifact) quantifies the low marginal runtime/model cost for deterministic rescue under the chosen repair budget and edit operator design.

Statistical interpretation and limitations of inference. The confirmatory inferential outcome (p = 1.0) primarily reflects the paucity of discordant pairs under the shared\_initial estimand and chosen synthetic suite. Exact McNemar inference depends only on n10 and n01; with n10 = 1 and n01 = 0 the two‑sided exact p equals 1.0, indicating no statistically detectable difference under conventional two‑sided testing in this sample. The paired bootstrap CI (seed = 1729, B = 10000) provides a complementary bound for effect size: the absolute paired improvement is plausibly between 0 and 0.01 given observed data and the preregistered estimand. These inferential statements apply strictly to the shared\_initial conditioning (one LLM draw reused across conditions), the limited repair budget (<=1 constrained edit), the synthetic task distribution, and the declared generator (gpt-5-mini) and validator (deterministic\_netops\_validator\_v5).

Practical diagnostic takeaway. Under operational policies that restrict model calls per change request, the guarded generate\ensuremath{\rightarrow}validate\ensuremath{\rightarrow}repair pattern produced a deterministic rescue in one hard example with negligible additional model‑call expenditure in this synthetic suite. However, because baseline single‑shot generation was near ceiling (299/300), the marginal observable benefit of bounded deterministic repair was correspondingly small in absolute terms. The archived per‑episode validator traces and repair transcript enable forensic inspection to inform whether richer repair operators, iterative repair budgets, or alternative estimand conditioning (e.g., allowing resampling) would be expected to materially increase rescue rates in different deployment regimes.

\section{Discussion}
Estimand conditioning and operational alignment. Conditioning on a single LLM draw per intent (shared\_initial) maps directly to operational policies that limit model calls for latency, cost, or governance reasons. Under that constraint the guarded pipeline measures deterministic rescue applied to the same candidate rather than gains from resampling or interactive steering. Consequently, the observed guarded benefit quantifies only the verifier27s and bounded repair operator27s ability to deterministically correct flaws present in an otherwise single generator output; it does not measure benefits that would accrue if the pipeline were allowed to request new independent candidates or to engage in multi‑turn verifier‑guided generation.

Stratum‑level signal and where deterministic rescue appears. Per‑stratum results (archived) show perfect baseline/guarded success in the low and medium strata (100/100 each) and a single baseline failure in the high stratum that the bounded repair corrected (baseline 99/100, guarded 100/100). This pattern indicates that, on this synthetic suite, generator single‑shot quality was largely sufficient for low/medium difficulty tasks and that deterministic rescue potential concentrated in rarer, higher‑difficulty sequencing omissions. The single discordant case diagnostic (a must\_be\_down\_for\_fields sequencing omission) illustrates a concrete failure class where a constrained edit can be effective: localized sequencing fixes that do not require reinterpretation of intent or broad structural changes.

Ceiling effects, pilot omission, and power implications. The preregistration targeted detection of a larger absolute improvement but relied on a pilot to estimate discordance; that pilot was not executed and the main run encountered a near‑ceiling baseline (299/300). As a result, discordant pairs were scarce (m = 1) and statistical sensitivity to small absolute improvements was necessarily limited. This outcome highlights an empirical lesson supported by the run27s logs: when generator single‑shot performance may approach ceiling on the chosen task suite, a small calibration pilot is valuable to estimate discordance and adjust sample sizing or task difficulty sampling before a full confirmatory run.

Validator coverage and failure‑mode dependence. The practical marginal utility of verifier‑assisted pipelines depends strongly on which failure modes occur in deployment and whether the validator expresses those properties. Our validator (deterministic\_netops\_validator\_v5) enforces sequencing invariants, required\_command\_count, preservation of unspecified state, and a set of topology reachability invariants encoded in the synthetic suite. On this suite most failure cases were absent or syntactic, producing few opportunities for deterministic rescue beyond localized edits. In operational contexts where failures more often involve richer reachability violations or ACL/policy semantics, a validator with broader semantic checks would detect additional failures and---if paired repair operators can safely correct them---could increase deterministic rescue rates. Conversely, if operator policy forbids automated edits, the guarded pipeline27s value is primarily as a precise detector and human‑in‑the‑loop escalation point rather than an autonomous repair engine.

Repair budget, cost, and auditability tradeoffs. We intentionally limited automated repair to a single constrained edit to bound potential operational risk and to measure low‑cost rescue opportunities. Execution accounting shows model\_calls\_used = 301 and repair\_stage\_count = 1, indicating that the observed deterministic rescue incurred one extra repair call for the single rescued episode. This low marginal model cost for the observed rescue suggests that, where localized sequencing omissions are the prevalent failure mode, a bounded repair policy can provide measurable benefit at modest additional model cost and with a compact repair trace for auditing. However, larger or iterative repair strategies, while potentially more effective, change the estimand and raise additional questions about edit semantics, cascading changes, and audit complexity.

Tradeoffs between validator complexity and generator resampling. Two distinct intervention families can improve pre‑deployment pass rates: (1) richer validators and constrained repairs that deterministically correct detectable defects in a given candidate, and (2) policies that permit resampling or verifier‑guided regeneration to exploit generator variance. These interventions have different operational costs and evidentiary interpretations. Rich validators require engineering effort to formalize semantics and maintain correctness; resampling increases model calls and may raise cost or governance concerns but can often yield higher empirical success rates. Evaluations should therefore align the experimental estimand with the intended operational policy: the shared\_initial estimand appropriately bounds deterministic rescue when resampling is disallowed, while alternative estimands are needed to evaluate resampling or interactive pipelines.

Practical implications for NetOps CI/CD and evaluation protocol design. For organizations that tightly budget model calls, our results provide an empirical upper bound on deterministic rescue achievable with a single shared\_initial candidate and a bounded (<=1) repair policy on the supplied synthetic suite. For organizations that permit resampling, evaluations should explicitly allow alternative candidate draws per condition. In either case, small calibration pilots to estimate discordance and stratified sampling across operationally relevant failure modes (sequencing, reachability, ACL/policy edits) will improve design sensitivity and help align evaluation outcomes with deployment decisions.

Synthesis with prior literature. The run27s pattern---limited deterministic rescue where generator single‑shot quality is high, with rescue concentrated in localized sequencing omissions---aligns with broader reports that tool‑augmentation and generate\ensuremath{\rightarrow}validate\ensuremath{\rightarrow}repair architectures can improve correctness but that realized gains depend on validator expressiveness and task distribution. Our artifact archive (validator traces and the single repair trace) provides concrete, inspectable evidence of the narrow failure class corrected here and can inform validator and repair operator design choices in future work.

Summary. The guarded pipeline27s deterministic rescue was real but marginal under the shared\_initial estimand on this synthetic suite and repair budget: one sequencing omission was fixed at low additional model cost. More substantive rescue would likely require (a) a task distribution with more semantic/structural failures detectable by the validator, (b) richer validator expressiveness (e.g., deeper reachability or ACL semantics), or (c) relaxing the estimand to allow resampling or interactive verifier‑guided generation. These conclusions are explicitly conditioned on the archived experimental design, artifacts, and the declared model/validator combination.

\section{Limitations}
\begin{itemize}
\item Synthetic benchmark: tasks and topologies were generated by the preregistered deterministic task generator; real‑world heterogeneity may expose different failure modes and increase verifier marginal utility.
\item Single model and validator: only gpt-5-mini and deterministic\_netops\_validator\_v5 were evaluated; results may not generalize to other LLMs or richer/diverse validators.
\item Ceiling effect: baseline single‑shot performance was near 100\% on this suite (299/300), producing limited discordance and reduced power to detect small absolute improvements relative to larger pre‑specified targets.
\item Pilot not executed: the preregistered pilot intended for discordance estimation and power calibration was not run; no post‑preregistration power recalculation was performed.
\item Estimand conditioning: the shared\_initial protocol conditions the estimand on a single LLM candidate and therefore does not measure verifier benefit that would arise from allowing independent alternative LLM candidates per condition (resampling or iterative generation).
\item Repair budget: the guarded condition permitted at most one automated repair per pair; larger or iterative repair strategies may yield different outcomes.
\end{itemize}

\section{Conclusion}
We executed a preregistered, artifact‑archived paired comparison between a verifier‑assisted pipeline (deterministic validator + <=1 automated repair) and accepting a single‑shot LLM candidate under the shared\_initial estimand on 300 synthetic NetOps intents using gpt-5-mini and deterministic\_netops\_validator\_v5. Baseline single‑shot generation produced 299/300 unflagged verifier passes and the guarded pipeline 300/300, yielding an absolute paired improvement of 0.00333 (exact McNemar two‑sided p = 1.0; 95\% bootstrap CI [0.00,0.01]). Deterministic validator+repair rescued a single sequencing omission under the bounded repair budget. The near‑ceiling baseline limits inferential sensitivity to small absolute improvements under the shared\_initial estimand. We release full preregistration, execution, and analysis artifacts to support reproducibility and targeted follow‑up studies that vary estimand conditioning, model strength, task distribution, or repair budget.

\section*{Disclosure Statement}
This experiment and manuscript were produced by an autonomous CNSM Agentic‑AI research run executed under the archived initial master prompt (provenance/master\_prompt.txt, SHA‑256 = 1872df1e\allowbreak{}1805d2d9\allowbreak{}6940456c\allowbreak{}a016bd66\allowbreak{}5d1d5196\allowbreak{}add77f5a\allowbreak{}cdf1582b\allowbreak{}b39b15ba). Human role: the human Research Director selected the topic family (Generative AI and LLMs for NetOps) and approved pre‑lock infrastructure; all literature search after the autonomy lock, autonomous study selection, preregistration (study id cnsmp2026-llm-verifier-predeployment-001), execution, analysis, manuscript drafting, AI peer review, and final compilation were performed autonomously under the locked master prompt. Models and orchestration: the declared generation model used was gpt-5-mini (execution/model\_configuration.json declares requested\_model = "gpt-5-mini"); adapter/orchestration framework: hosted\_netops\_gvr\_v1. Validator and tooling: deterministic\_netops\_validator\_v5 performed semantic and syntactic validation according to the preregistered full\_intent\_constraint\_validation rule; the guarded condition permitted at most one automated repair call per pair implemented as a constrained edit operator. Preregistration: the confirmatory design, estimand, analysis executor, exclusion/missingness rules, and multiplicity plan were frozen prior to the main run and archived with the study id (preregistration SHA‑256 = 425cb165\allowbreak{}2354120c\allowbreak{}f48f686b\allowbreak{}964723d6\allowbreak{}4c76d0a8\allowbreak{}abe53cd5\allowbreak{}b5ad0eaa\allowbreak{}1af71fc7). Execution summary (archived): planned\_episode\_count = 600, completed\_episode\_count = 600, complete\_pair\_count = 300, model\_calls\_used = 301; provider audit: hosted\_model\_call\_event\_count = 301, stage\_counts show shared\_initial\_generation = 300 and repair = 1. Pilot: the preregistration described an optional pilot (n = 60) for discordance estimation and power calibration; that pilot was not executed and no post‑preregistration recalculation of required sample size was performed. Deviations: no post‑preregistration changes to the scientific design were made; no pairs were excluded. Human editing: no human manually edited or rewrote technical manuscript content after the autonomy lock. Artifact availability: all raw outputs, logs, task manifests, validator traces, repair provider traces, and analysis artifacts referenced in this manuscript are archived in the run bundle associated with the study id. The archived run bundle contains the low‑level repair provider trace documenting the single automated repair and subsequent validation pass; this manuscript reports a concise semantic summary of that repair in Results rather than reproducing full verbatim provider transcripts in the body.

\begin{thebibliography}{99}
\bibitem{ref1} Yajie Zhou, Kevin Hsieh, Sathiya Kumaran Mani, Srikanth Kandula, Zaoxing Liu, "MeshAgent: Enabling Reliable Network Management with Large Language Models", 2025, 10.1145/3771567
\bibitem{ref2} Xiaolong Wei, Yuehu Dong, Xingliang Wang, Xingyu Zhang, Zhejun Zhao, Dongdong Shen, Long Xia, Dawei Yin, "Beyond ReAct: A Planner-Centric Framework for Complex Tool-Augmented LLM Reasoning", 2026, 10.1609/aaai.v40i40.40676
\bibitem{ref3} Xu Liu, Peng Zhang, Anubhavnidhi Abhashkumar, Jiawei Chen, Weirong Jiang, "Automatic Configuration Repair", 2024, 10.1145/3696348.3696895
\bibitem{ref4} Hao Zhou, Chengming Hu, Ye Yuan, Yufei Cui, Yili Jin, Can Chen, Haolun Wu, Dun Yuan, Li Jiang, Di Wu, Xue Liu, Jianzhong Zhang, Xianbin Wang, Jiangchuan Liu, "Large Language Model (LLM) for Telecommunications: A Comprehensive Survey on Principles, Key Techniques, and Opportunities", 2024, 10.1109/comst.2024.3465447
\bibitem{ref5} X X Zhang, Xianming Gao, Peilin Tao, Tao Feng, "GraphSynth: Synthesis of Network Configuration Templates Using Large Language Models", 2025, 10.1145/3728725.3728742
\bibitem{ref6} Sebastian Farquhar, Jannik Kossen, Lorenz Kuhn, Yarin Gal, "Detecting hallucinations in large language models using semantic entropy", 2024, 10.1038/s41586-024-07421-0
\bibitem{ref7} Yunze Wei, Wei, Kaiwen, Shaohui Du, Wang, Jianyu, Liu, Zhangzhong, Yawen Wang, Zhenxin Li, Congcong Miao, Xiaohui Xie, Yong Cui, "Automated Network Protocol Testing with LLM Agents", 2025, 10.48550/arxiv.2510.13248
\bibitem{ref8} Fuliang Li, Haozhi Lang, Jiajie Zhang, Jiaxing Shen, Xingwei Wang, "PreConfig: A Pretrained Model for Automating Network Configuration", 2024, 10.48550/arxiv.2403.09369
\end{thebibliography}

\end{document}
